\documentclass[conference]{IEEEtran}

\usepackage{cite}
\usepackage{amsmath,amssymb}
\usepackage{url}
\usepackage{listings}
\usepackage{xcolor}

\lstdefinestyle{ocamlstyle}{
  basicstyle=\ttfamily\footnotesize,
  breaklines=true,
  columns=fullflexible,
  frame=single,
  numbers=left,
  numberstyle=\tiny,
  keywordstyle=\color{blue},
  stringstyle=\color{teal}
  commentstyle=\color{gray},
}

\title{Mid-Evaluation Report: Functional Planner Kernels for Dagents}

\author{
\IEEEauthorblockN{Pradyun Devarakonda \quad Dheeraj Rajanala}
\IEEEauthorblockA{
Course Project for Programming Languages\\
Roll Number: IMT2022525, IMT2022093}
}

\begin{document}
\maketitle

% \begin{abstract}
% This report presents the mid-evaluation status of the functional programming
% component of Dagents, a reusable agentic framework for machine learning
% automation, data engineering, model routing, pipeline planning, and workload
% manifest generation. The functional programming side of the project focuses on
% using OCaml as a typed planner layer rather than as a replacement
% for Python or Spring Boot services. The current work implements deterministic
% functional kernels for shared intermediate representations, dataset schema
% planning, source validation, extraction planning, data quality checks, pipeline
% DAG planning, model routing, and Kubernetes manifest generation. The main
% programming languages objective is to evaluate how algebraic data types,
% pattern matching, immutability, and pure transformations can improve reliability
% in data engineering and ML orchestration workflows.
% \end{abstract}

\begin{IEEEkeywords}
functional programming, OCaml, domain-specific planning, data engineering,
pipeline planning, type systems, ML orchestration, Kubernetes
\end{IEEEkeywords}

\section{Problem Description}

Dagents is a shared framework intended to support reusable agentic services for
machine learning automation. Its broader architecture contains Local Monitoring
Agents (LMA), Global Monitoring Agents (GMA), model execution services, pipeline
services, and core orchestration services. The functional programming part of
this project addresses a specific problem inside that architecture: many
orchestration decisions are currently expressed as ad hoc service logic, even
though they are better understood as deterministic planning, translation, and
validation problems.

The core problem is to design a functional programming module that can translate
declarative specifications into execution-ready plans. Examples include
lowering a data source specification into an extraction plan, lowering a
pipeline definition into a topologically ordered execution graph, routing a
dataset profile to candidate model families, and rendering workload
descriptions as Kubernetes manifests. These tasks are not primarily concerned
with stateful I/O. They are transformations from one structured representation
to another, which makes them a natural fit for functional programming.

In this report, the term ``compiler'' is used only in the domain-specific sense:
it means a planner or translator from one Dagents representation to another.
It does not mean a programming-language compiler that parses a source language,
performs semantic analysis for a general-purpose language, or generates machine
code. The implemented OCaml modules are closer to typed validators, planners,
and renderers for data engineering workflows. \\

Repo Link: \url{https://github.com/pradyunuydarp/Dagents.git}

\subsection{Learning Objectives}

The project has the following learning objectives from the perspective of a
Programming Languages course:

\begin{itemize}
  \item Study how algebraic data types can model domain concepts such as source
  kinds, task types, pipeline steps, model families, data quality operators, and
  workload kinds.
  \item Use pattern matching to make planning decisions explicit and exhaustive.
  \item Apply immutability and pure functions to build reproducible planning
  logic for data engineering and ML workflows.
  \item Separate pure planning logic from runtime effects such as database
  access, model training, API routing, and Kubernetes deployment.
  \item Evaluate how a typed functional kernel can coexist with Python and Java
  services through process-level JSON contracts.
\end{itemize}

\subsection{Intended Outcomes}

By the end of the project, the functional programming side should provide:

\begin{itemize}
  \item A documented OCaml codebase containing reusable functional planner modules.
  \item A common intermediate representation (IR) for datasets, sources,
  pipelines, models, and workloads.
  \item A dataset planner for schema inference, source validation, extraction
  planning, quality rule evaluation, and transform planning.
  \item A pipeline planner for dependency validation, cycle detection, and
  execution target assignment.
  \item A model router that maps dataset profiles to model families and
  packaging strategies.
  \item A manifest renderer that generates deterministic Kubernetes YAML from
  typed workload specifications.
  \item A command-line boundary, \texttt{dagentsc}, that exposes the OCaml
  functional kernels to other services without requiring direct FFI.
\end{itemize}

\section{Solution Outline}

The solution is structured around the idea that OCaml should be used for
functional kernels rather than for the entire distributed system. Python remains
appropriate for numerical model training and data science libraries. Spring
Boot remains appropriate for API orchestration and external integration.
OCaml is used where the project benefits from typed, deterministic, and
planner-like transformations.

\subsection{Architecture}

The current OCaml workspace is organized as follows:

\begin{lstlisting}[style=ocamlstyle]
bindings/ocaml
|-- lib/common_ir
|-- lib/dataset_compiler
|-- lib/pipeline_compiler
|-- lib/model_router
|-- lib/manifest_compiler
|-- bin/dagentsc.ml
`-- test/test_compilers.ml
\end{lstlisting}

The \texttt{common\_ir} module defines shared types. These include source
specifications, schema contracts, dataset profiles, quality rules, transform
operations, pipeline definitions, model route plans, and workload manifests.
The individual planner modules consume these typed structures and emit
normalized plans. The directory names use \texttt{compiler} because the modules
perform domain-specific lowering, but the work is not about implementing a
general-purpose programming-language compiler.

\subsection{Common Intermediate Representation}

The IR is the foundation of the functional design. Instead of passing loosely
structured maps between modules, the project defines explicit types such as:

\begin{itemize}
  \item \texttt{source\_kind}: inline, Postgres, MongoDB, or object storage.
  \item \texttt{task\_type}: anomaly detection, classification, forecasting,
  embedding, or regression.
  \item \texttt{step\_kind}: enrichment, filtering, projection, profiling, or
  model execution.
  \item \texttt{quality\_operator}: non-null, unique, minimum value, maximum
  value, glob-style match, or allowed values.
  \item \texttt{workload\_kind}: Deployment, Job, CronJob, Service, or
  ConfigMap.
\end{itemize}

This is important for programming language design because domain rules are
encoded in the type structure. When new source kinds or pipeline steps are
added, pattern matches in the planner modules reveal the places that must be
updated.

\subsection{Dataset Planner}

The dataset planner is the main data engineering module. It currently
supports:

\begin{itemize}
  \item \textbf{Schema inference}: derive stable field names and scalar data
  types from inline records.
  \item \textbf{Source validation}: reject malformed source specifications,
  such as external sources without connection references or Postgres sources
  without either SQL or table selection.
  \item \textbf{Extraction planning}: lower source-specific selections into a
  normalized extraction plan containing selected fields, predicates, ordering,
  batch size, checkpoint, and partition strategy.
  \item \textbf{Schema contract validation}: compare actual schemas against
  required and optional fields.
  \item \textbf{Quality rule evaluation}: count violations for rules such as
  non-null, uniqueness, numeric bounds, pattern match, and allowed values.
  \item \textbf{Transform planning}: lower declarative operations such as
  select, drop, rename, cast, and non-null filtering into an output schema.
  \item \textbf{Transform execution}: apply simple record-level transformations
  for local or test-time data engineering workflows.
\end{itemize}

A key design decision is that the dataset planner does not open database
connections or object storage readers. It only validates and plans. Runtime
adapters in Python, LMA, GMA, or pipeline services perform actual I/O. This
preserves functional purity and keeps failure isolation clear.

\subsection{Pipeline Planner}

The pipeline planner treats a workflow as a directed acyclic graph (DAG). It
performs the following checks and transformations:

\begin{itemize}
  \item reject duplicate step identifiers;
  \item reject unknown dependencies;
  \item detect cyclic dependencies;
  \item produce a topologically ordered step list;
  \item assign execution targets such as local process or Python service.
\end{itemize}

This module shows planner-style use of functional programming: the input is a
high-level declarative pipeline, and the output is a lower-level execution plan.

\subsection{Model Router}

The model router maps a \texttt{dataset\_profile} and a \texttt{task\_type} to
a route plan. The route plan contains candidate model families, the selected
default model, and a packaging mode. For example, forecasting routes toward
sequence models and long-running deployment, while tabular anomaly detection can
route toward autoencoders or Kubernetes job execution for large datasets.

\subsection{Manifest Renderer}

The manifest renderer converts typed workload specifications into deterministic
Kubernetes YAML. It supports Deployments, Jobs, CronJobs, Services, and
ConfigMaps. The renderer applies defaults for resources, ports, environment
variables, and labels. This module demonstrates how functional programming can
be used for infrastructure synthesis: a structured workload specification is
rendered into deployment artifacts.

\subsection{Command-Line Boundary}

The \texttt{dagentsc} executable exposes the OCaml modules over a process
boundary. Example commands include:

\begin{lstlisting}[style=ocamlstyle]
dagentsc dataset source validate --input source.json
dagentsc dataset source extract --input source.json
dagentsc dataset schema validate --records records.json --contract contract.json
dagentsc dataset quality evaluate --records records.json --rules rules.json
dagentsc dataset transform apply --records records.json --operations operations.json
dagentsc pipeline compile --input pipeline.json --output json
dagentsc manifest compile --input workload.json --output json
\end{lstlisting}

This avoids direct runtime embedding of OCaml into Python or Java. It also makes
the contract explicit and testable.

\section{Expected Deliverables}

The final deliverables are expected to include:

\begin{itemize}
  \item \textbf{Codebase}: OCaml modules under \texttt{bindings/ocaml} for IR,
  dataset planning, pipeline planning, model routing, manifest rendering, JSON
  codecs, CLI integration, and tests.
  \item \textbf{Study report}: explanation of why functional programming is
  suitable for deterministic planning in Dagents.
  \item \textbf{Algorithms}: DAG validation and ordering, schema inference,
  source validation, extraction plan generation, quality rule evaluation, and
  transform planning.
  \item \textbf{Conceptual framework}: a plan-then-execute architecture for
  data engineering and ML orchestration.
  \item \textbf{Documentation}: module design notes, CLI examples, extension
  guidelines, and report material for evaluation.
\end{itemize}

\section{Main References}

The project uses the following key references. Tutorials and videos are
intentionally excluded.

\begin{enumerate}
  \item \textbf{Milner, ``A Theory of Type Polymorphism in Programming,'' 1978}
  \cite{milner1978}. This paper is relevant because the project depends on the
  idea that static type systems can express generic and reusable program
  structure. The OCaml kernels use types to represent reusable planner
  contracts.

  \item \textbf{Leroy et al., ``The OCaml System: Documentation and User's
  Manual''} \cite{ocamlmanual}. This is the primary tool reference for the
  language used in the functional module. It supports the project's use of
  algebraic data types, modules, pattern matching, and compiled native
  executables.

  \item \textbf{Dean and Ghemawat, ``MapReduce: Simplified Data Processing on
  Large Clusters,'' 2004} \cite{mapreduce}. This paper motivates the data
  engineering side of the project. Although Dagents is not implementing
  MapReduce directly, the paper frames large-scale data processing as planned,
  partitioned, deterministic transformations.

  \item \textbf{Kubernetes Documentation, ``Workloads''} \cite{kubernetes}.
  Kubernetes is relevant because Dagents renders typed workload specifications
  into Deployment, Job, CronJob, Service, and ConfigMap manifests. This
  reference defines the target infrastructure concepts generated by the
  manifest renderer.
\end{enumerate}

\section{Team Progress}

\subsection{Completed Milestones}

The team has completed the following milestones by mid-evaluation:

\begin{itemize}
  \item Established an OCaml workspace using Dune.
  \item Implemented a shared \texttt{common\_ir} module containing typed
  contracts for source specs, dataset profiles, pipelines, models, quality
  rules, transforms, and workload manifests.
  \item Implemented JSON codecs for the command-line and service boundary.
  \item Implemented source validation and extraction planning for inline,
  Postgres, MongoDB, and object storage source types.
  \item Implemented schema inference and schema contract validation.
  \item Implemented quality rule evaluation for common data engineering checks.
  \item Implemented transform planning and local transform application.
  \item Implemented pipeline DAG validation, cycle detection, and topological
  ordering.
  \item Implemented model routing from dataset profiles to model families and
  packaging modes.
  \item Implemented Kubernetes manifest generation for common workload kinds.
  \item Added a process-boundary CLI, \texttt{dagentsc}, for service
  integration.
  \item Added OCaml tests covering planner behavior and verified them using
  \texttt{opam exec -- dune test}.
  \item Added design documentation for the OCaml functional modules.
\end{itemize}

\subsection{Achievements}

The main achievement is that the functional programming module is no longer
only a conceptual plan. It is now a working codebase with typed contracts,
planner functions, CLI integration, and tests. The system demonstrates how
functional programming can simplify orchestration logic by moving decision
rules into small deterministic modules.

Another achievement is the separation of pure logic from effects. For example,
the source planner validates and plans a Postgres extraction, but it does not
connect to Postgres. This separation makes the OCaml layer easier to test and
allows runtime services to evolve independently.

\subsection{Challenges}

Several challenges were encountered:

\begin{itemize}
  \item Designing IR types that are general enough for future services but not
  so abstract that they become hard to test.
  \item Keeping JSON contracts compatible with Python and Java service
  boundaries while still preserving typed OCaml internals.
  \item Deciding which data engineering tasks should remain in OCaml and which
  should remain in Python runtime adapters.
  \item Balancing report-level conceptual clarity with practical code
  implementation.
  \item Avoiding direct FFI complexity by using a process boundary while still
  making the CLI ergonomic.
\end{itemize}

\section{Work Distribution}

There are two team members in the project.

\subsection{Dheeraj: Functional Core and Planner Design}

Dheeraj is responsible for the functional programming core:

\begin{itemize}
  \item designing the OCaml intermediate representation;
  \item implementing algebraic data types for sources, schemas, pipelines,
  model families, quality rules, transforms, and workload specs;
  \item implementing the dataset planner;
  \item implementing schema inference, schema validation, quality evaluation,
  and transform planning;
  \item writing OCaml unit tests for planner correctness;
  \item ensuring that the code follows functional programming principles such
  as immutability, pattern matching, and pure transformations.
\end{itemize}

\subsection{Pradyun: Integration, Service Boundary, and Documentation}

Pradyun is responsible for integration and project framing:

\begin{itemize}
  \item designing the \texttt{dagentsc} command-line interface;
  \item implementing JSON input and output contracts for service integration;
  \item connecting the OCaml planner concept to the broader Dagents
  architecture involving LMA, GMA, pipeline service, model service, and core
  service;
  \item documenting the module responsibilities, CLI examples, extension
  guidelines, and architecture decisions;
  \item validating the manifest renderer and process-boundary integration
  strategy;
  \item writing OCaml grey box and black box tests for planner correctness;
  \item preparing report material and aligning the project with Programming
  Languages course objectives.
\end{itemize}

\section{Next Steps}

The remaining work focuses on hardening the module and improving the evaluation
story:

\begin{itemize}
  \item Add more negative tests for invalid source specs, invalid schemas, and
  cyclic pipelines.
  \item Expand quality rules with severity-aware blocking decisions.
  \item Add richer partition planning, especially time-window partitioning.
  \item Connect more service calls to \texttt{dagentsc} in a controlled way.
  \item Compare equivalent planning logic in Python and OCaml to evaluate
  readability, safety, and maintainability.
  \item Finalize the end-term report with examples, test evidence, and
  discussion of programming language tradeoffs.
\end{itemize}

\section{Conclusion}

At mid-evaluation, the functional programming side of Dagents has produced a
working OCaml-based planner layer. The project uses functional
programming not as a general replacement for all services, but as a precise tool
for typed transformations. This makes the project suitable for a Programming
Languages course because it demonstrates practical use of algebraic data types,
pattern matching, pure functions, domain-specific lowering, and process-boundary
language integration in a realistic data engineering and ML automation system.

\begin{thebibliography}{1}

\bibitem{milner1978}
R. Milner, ``A theory of type polymorphism in programming,'' \emph{Journal of
Computer and System Sciences}, vol. 17, no. 3, pp. 348--375, 1978.

\bibitem{ocamlmanual}
X. Leroy, D. Doligez, A. Frisch, J. Garrigue, D. Remy, and J. Vouillon,
\emph{The OCaml System: Documentation and User's Manual}. OCaml.org.
[Online]. Available: \url{https://ocaml.org/manual/}

\bibitem{mapreduce}
J. Dean and S. Ghemawat, ``MapReduce: Simplified data processing on large
clusters,'' in \emph{Proc. 6th USENIX Symposium on Operating Systems Design and
Implementation (OSDI)}, 2004, pp. 137--150.

\bibitem{kubernetes}
The Kubernetes Authors, ``Workloads,'' Kubernetes Documentation. [Online].
Available: \url{https://kubernetes.io/docs/concepts/workloads/}

\end{thebibliography}

\end{document}
